\documentclass[11pt,letterpaper]{article}

\usepackage[margin=2.2cm]{geometry}
\usepackage{amsmath}
\usepackage{amssymb}
\usepackage{array}
\usepackage{booktabs}

\title{Memoria de calculo detallada\\Proyecto IE0413}
\author{Grupo: completar nombres y carnes}
\date{I ciclo 2026}

\begin{document}

\maketitle

\section*{Nota de uso}

Este archivo es una base para la memoria de calculo. La realimentacion recibida
indica que los calculos pueden escribirse en \LaTeX, siempre que se muestre el
desarrollo paso a paso. Por esa razon, en cada etapa se escriben las ecuaciones,
los despejes y la sustitucion numerica.

Los amplificadores operacionales usados en la simulacion son LF353 no ideales,
mediante el subcircuito \texttt{LF353/NS}. Los calculos manuales usan las
relaciones de realimentacion vistas en clase y luego se verifican con el
macromodelo del LF353 en LTspice.

\section{Parametros definidos por los carnes}

Primero se ordenan alfabeticamente los primeros apellidos de los tres
integrantes. Luego se toma el primer digito no nulo del carne de cada uno:

\[
X = \text{digito del integrante 1}, \qquad
Y = \text{digito del integrante 2}, \qquad
Z = \text{digito del integrante 3}.
\]

La frecuencia pedida para el oscilador central es:

\[
f_0 = X,YZ~\text{kHz}.
\]

Para escribirla como una expresion decimal:

\[
f_0 =
\left(
X + \frac{Y}{10} + \frac{Z}{100}
\right) 1000~\text{Hz}.
\]

Ejemplo usado en los archivos de simulacion:

\[
X=1,\qquad Y=2,\qquad Z=3.
\]

Entonces:

\[
f_0 =
\left(
1+\frac{2}{10}+\frac{3}{100}
\right)1000
=
\left(
1+0.2+0.03
\right)1000
=1.23(1000)
=1230~\text{Hz}.
\]

Por lo tanto:

\[
f_0 = 1.23~\text{kHz}.
\]

\section{Oscilador central de onda cuadrada}

\subsection{Topologia}

Se usa un oscilador de relajacion formado por:

\begin{itemize}
  \item un comparador Schmitt con LF353,
  \item una red resistiva de realimentacion positiva,
  \item un capacitor que se carga y descarga por medio de una resistencia.
\end{itemize}

La salida del LF353 satura aproximadamente en dos niveles simetricos. Para el
calculo se escriben como:

\[
v_o = +V_S
\qquad \text{y} \qquad
v_o = -V_S.
\]

El divisor de realimentacion fija el voltaje de referencia del disparador
Schmitt:

\[
v_+ = \beta v_o.
\]

Si el divisor esta formado por dos resistencias, una desde la salida hasta el
nodo no inversor y otra desde ese nodo hasta tierra, entonces:

\[
\beta = \frac{R_2}{R_1+R_2}.
\]

Se escogio:

\[
R_1 = R_2 = 100~\text{k}\Omega.
\]

Sustituyendo:

\[
\beta =
\frac{100~\text{k}\Omega}{100~\text{k}\Omega+100~\text{k}\Omega}
=
\frac{100}{200}
=0.5.
\]

Los umbrales del Schmitt son:

\[
V_{TH+}=+\beta V_S,
\qquad
V_{TH-}=-\beta V_S.
\]

Con \(\beta=0.5\):

\[
V_{TH+}=+0.5V_S,
\qquad
V_{TH-}=-0.5V_S.
\]

\subsection{Deduccion del periodo}

Durante un semiciclo, suponga que la salida esta en \(+V_S\). El capacitor se
carga hacia \(+V_S\) desde el valor inicial \(-\beta V_S\). La ecuacion de carga
del capacitor es:

\[
v_C(t)=V_S+\left(v_C(0)-V_S\right)e^{-t/(RC)}.
\]

Al inicio de ese semiciclo:

\[
v_C(0)=-\beta V_S.
\]

Sustituyendo:

\[
v_C(t)=V_S+\left(-\beta V_S-V_S\right)e^{-t/(RC)}.
\]

Se factoriza \(V_S\):

\[
v_C(t)=V_S-(1+\beta)V_S e^{-t/(RC)}.
\]

El cambio de estado ocurre cuando el capacitor alcanza el umbral positivo:

\[
v_C(t_h)=+\beta V_S.
\]

Entonces:

\[
\beta V_S = V_S-(1+\beta)V_S e^{-t_h/(RC)}.
\]

Se divide toda la ecuacion entre \(V_S\):

\[
\beta = 1-(1+\beta)e^{-t_h/(RC)}.
\]

Se despeja el termino exponencial:

\[
(1+\beta)e^{-t_h/(RC)}=1-\beta.
\]

\[
e^{-t_h/(RC)}=\frac{1-\beta}{1+\beta}.
\]

Se aplica logaritmo natural a ambos lados:

\[
-\frac{t_h}{RC}
=
\ln \left(
\frac{1-\beta}{1+\beta}
\right).
\]

Multiplicando por \(-RC\):

\[
t_h =
RC
\ln \left(
\frac{1+\beta}{1-\beta}
\right).
\]

Como la carga y la descarga tienen la misma duracion:

\[
T = 2t_h.
\]

Por tanto:

\[
T =
2RC
\ln \left(
\frac{1+\beta}{1-\beta}
\right).
\]

La frecuencia es:

\[
f_0=\frac{1}{T}.
\]

Entonces:

\[
f_0 =
\frac{1}{
2RC
\ln \left(
\frac{1+\beta}{1-\beta}
\right)
}.
\]

Se despeja \(R\):

\[
2f_0RC
\ln \left(
\frac{1+\beta}{1-\beta}
\right)
=1.
\]

\[
R =
\frac{1}{
2f_0C
\ln \left(
\frac{1+\beta}{1-\beta}
\right)
}.
\]

\subsection{Sustitucion numerica}

Se escogio:

\[
C_{osc}=10~\text{nF}=10\times 10^{-9}~\text{F}.
\]

Ademas:

\[
\beta=0.5.
\]

Entonces:

\[
\ln\left(
\frac{1+\beta}{1-\beta}
\right)
=
\ln\left(
\frac{1+0.5}{1-0.5}
\right)
=
\ln\left(
\frac{1.5}{0.5}
\right)
=
\ln(3).
\]

Para el ejemplo \(f_0=1230~\text{Hz}\):

\[
R_{osc,teo} =
\frac{1}{
2(1230)(10\times 10^{-9})\ln(3)
}.
\]

Primero:

\[
2(1230)(10\times 10^{-9})
=24600\times 10^{-9}
=2.46\times 10^{-5}.
\]

Luego:

\[
(2.46\times 10^{-5})\ln(3)
=(2.46\times 10^{-5})(1.0986)
=2.7026\times 10^{-5}.
\]

Finalmente:

\[
R_{osc,teo}
=
\frac{1}{2.7026\times 10^{-5}}
=37001.6~\Omega
=37.0~\text{k}\Omega.
\]

En LTspice se usa el macromodelo no ideal del LF353. Para que la frecuencia de
simulacion coincida con la frecuencia pedida, se aplica el factor de ajuste:

\[
K_{trim}=1.0462.
\]

El valor usado en simulacion es:

\[
R_{osc,LTspice}=K_{trim}R_{osc,teo}.
\]

Sustituyendo:

\[
R_{osc,LTspice}
=
(1.0462)(37001.6~\Omega)
=
38711.1~\Omega.
\]

Por tanto:

\[
R_{osc,LTspice}\approx 38.7~\text{k}\Omega.
\]

\section{Onda diente de sierra positiva}

\subsection{Frecuencia requerida}

El enunciado pide que la onda diente de sierra se repita en los semiciclos
positivo y negativo de la onda cuadrada. Por eso:

\[
f_{saw}=2f_0.
\]

Para el ejemplo:

\[
f_{saw}=2(1230)=2460~\text{Hz}.
\]

El periodo de la diente de sierra es:

\[
T_{saw}=\frac{1}{f_{saw}}=\frac{1}{2f_0}.
\]

Para el ejemplo:

\[
T_{saw}
=
\frac{1}{2460}
=
4.0650\times 10^{-4}~\text{s}
=0.4065~\text{ms}.
\]

\subsection{Relacion entre corriente, capacitor y pendiente}

Si un capacitor se carga con corriente constante:

\[
i_C=C\frac{dv_C}{dt}.
\]

Como la corriente es constante:

\[
\frac{dv_C}{dt}=\frac{I}{C}.
\]

La variacion de voltaje durante un periodo de rampa es:

\[
\Delta V = V_{saw,pk}-0=V_{saw,pk}.
\]

La pendiente tambien puede escribirse como:

\[
\frac{\Delta V}{\Delta t}
=
\frac{V_{saw,pk}}{T_{saw}}.
\]

Igualando ambas expresiones de pendiente:

\[
\frac{I_{saw}}{C_{saw}}
=
\frac{V_{saw,pk}}{T_{saw}}.
\]

Se despeja la corriente:

\[
I_{saw}
=
C_{saw}
\frac{V_{saw,pk}}{T_{saw}}.
\]

Como:

\[
T_{saw}=\frac{1}{2f_0},
\]

entonces:

\[
I_{saw}
=
C_{saw}V_{saw,pk}(2f_0).
\]

\subsection{Sustitucion numerica}

Se escogio:

\[
C_{saw}=10~\text{nF}=10\times 10^{-9}~\text{F},
\qquad
V_{saw,pk}=5~\text{V}.
\]

Entonces:

\[
I_{saw}
=
(10\times 10^{-9})(5)(2)(1230).
\]

Primero:

\[
(5)(2)(1230)=12300.
\]

Luego:

\[
I_{saw}
=
(10\times 10^{-9})(12300)
=
123000\times 10^{-9}
=
123\times 10^{-6}~\text{A}.
\]

Por tanto:

\[
I_{saw}=123~\mu\text{A}.
\]

\section{Onda triangular}

\subsection{Frecuencia requerida}

El enunciado indica que la frecuencia de salida triangular es la mitad de la del
oscilador central:

\[
f_{tri}=\frac{f_0}{2}.
\]

Para el ejemplo:

\[
f_{tri}
=
\frac{1230}{2}
=
615~\text{Hz}.
\]

El periodo triangular es:

\[
T_{tri}
=
\frac{1}{f_{tri}}
=
\frac{1}{615}
=
1.6260\times 10^{-3}~\text{s}
=
1.626~\text{ms}.
\]

\subsection{Cambiador de nivel}

La salida equivalente del divisor por dos se toma entre \(0~\text{V}\) y
\(2~\text{V}\). Para usarla en el integrador con excursion positiva y negativa,
se resta \(1~\text{V}\):

\[
v_{drv}=v_{div}-1~\text{V}.
\]

Entonces, si:

\[
v_{div}=0~\text{V},
\]

se obtiene:

\[
v_{drv}=0-1=-1~\text{V}.
\]

Y si:

\[
v_{div}=2~\text{V},
\]

se obtiene:

\[
v_{drv}=2-1=+1~\text{V}.
\]

Por tanto, el integrador recibe una onda cuadrada centrada en cero:

\[
v_{drv}=\pm 1~\text{V}.
\]

\subsection{Integrador inversor}

Para un integrador inversor con LF353, resistencia de entrada \(R_{tri}\) y
capacitor de realimentacion \(C_{tri}\):

\[
i_R=\frac{v_{drv}}{R_{tri}}.
\]

Como la entrada inversora esta en el nodo de suma, la corriente de la resistencia
pasa por el capacitor:

\[
i_C=i_R.
\]

En el capacitor:

\[
i_C=C_{tri}\frac{d(v_- - v_o)}{dt}.
\]

Como el nodo de suma se mantiene cercano a cero por realimentacion:

\[
v_- \approx 0.
\]

Entonces:

\[
i_C=C_{tri}\frac{d(0-v_o)}{dt}
=
-C_{tri}\frac{dv_o}{dt}.
\]

Igualando corrientes:

\[
\frac{v_{drv}}{R_{tri}}
=
-C_{tri}\frac{dv_o}{dt}.
\]

Se despeja la pendiente:

\[
\frac{dv_o}{dt}
=
-\frac{v_{drv}}{R_{tri}C_{tri}}.
\]

La magnitud de la pendiente es:

\[
\left|\frac{dv_o}{dt}\right|
=
\frac{V_{drv}}{R_{tri}C_{tri}}.
\]

La onda triangular cambia de \(-V_{tri,pk}\) a \(+V_{tri,pk}\) durante medio
periodo de la triangular. Como:

\[
f_{tri}=\frac{f_0}{2},
\]

el medio periodo triangular es:

\[
\frac{T_{tri}}{2}
=
\frac{1}{2f_{tri}}
=
\frac{1}{2(f_0/2)}
=
\frac{1}{f_0}.
\]

La variacion de voltaje durante ese intervalo es:

\[
\Delta V = (+V_{tri,pk})-(-V_{tri,pk})=2V_{tri,pk}.
\]

La magnitud de la pendiente tambien es:

\[
\left|\frac{dv_o}{dt}\right|
=
\frac{2V_{tri,pk}}{1/f_0}
=
2V_{tri,pk}f_0.
\]

Se igualan las dos expresiones de pendiente:

\[
\frac{V_{drv}}{R_{tri}C_{tri}}
=
2V_{tri,pk}f_0.
\]

Se despeja \(R_{tri}\):

\[
R_{tri}
=
\frac{V_{drv}}{2V_{tri,pk}f_0C_{tri}}.
\]

\subsection{Sustitucion numerica}

Se escogio:

\[
V_{drv}=1~\text{V},
\qquad
V_{tri,pk}=2~\text{V},
\qquad
C_{tri}=10~\text{nF}.
\]

Entonces:

\[
R_{tri}
=
\frac{1}{
2(2)(1230)(10\times 10^{-9})
}.
\]

Primero:

\[
2(2)(1230)=4920.
\]

Luego:

\[
4920(10\times 10^{-9})
=
49200\times 10^{-9}
=
4.92\times 10^{-5}.
\]

Finalmente:

\[
R_{tri}
=
\frac{1}{4.92\times 10^{-5}}
=
20325.2~\Omega.
\]

Por tanto:

\[
R_{tri}\approx 20.3~\text{k}\Omega.
\]

\section{Seleccion de la onda para obtener la senoidal}

Se comparan las tres ondas disponibles desde el punto de vista de sus armonicos.

\subsection{Onda cuadrada}

Una onda cuadrada simetrica contiene solamente armonicos impares:

\[
v_{cuad}(t)=
\frac{4A}{\pi}
\left[
\sin(\omega_0 t)
+\frac{1}{3}\sin(3\omega_0 t)
+\frac{1}{5}\sin(5\omega_0 t)
+\cdots
\right].
\]

La amplitud del armonico \(n\) impar cae proporcionalmente a:

\[
\frac{1}{n}.
\]

Por ejemplo, el tercer armonico tiene amplitud relativa:

\[
\frac{1}{3}=0.333.
\]

\subsection{Onda diente de sierra}

Una onda diente de sierra contiene armonicos pares e impares:

\[
v_{saw}(t)=
\frac{2A}{\pi}
\left[
\sin(\omega_0t)
-\frac{1}{2}\sin(2\omega_0t)
+\frac{1}{3}\sin(3\omega_0t)
-\cdots
\right].
\]

La amplitud de los armonicos tambien cae proporcionalmente a:

\[
\frac{1}{n}.
\]

Ademas, al tener armonicos pares e impares, hay mas componentes que deben ser
atenuadas por el filtro pasabajas.

\subsection{Onda triangular}

Una onda triangular simetrica contiene solamente armonicos impares:

\[
v_{tri}(t)=
\frac{8A}{\pi^2}
\left[
\sin(\omega_0t)
-\frac{1}{3^2}\sin(3\omega_0t)
+\frac{1}{5^2}\sin(5\omega_0t)
-\cdots
\right].
\]

La amplitud del armonico \(n\) impar cae proporcionalmente a:

\[
\frac{1}{n^2}.
\]

Para el tercer armonico:

\[
\frac{1}{3^2}
=
\frac{1}{9}
=0.111.
\]

Comparando con la onda cuadrada:

\[
\frac{1/9}{1/3}
=
\frac{1}{3}.
\]

Es decir, antes de filtrar, el tercer armonico de la triangular ya es tres veces
menor que el de una cuadrada con la misma fundamental normalizada. Por esta
razon se escoge la onda triangular como entrada del filtro pasabajas de orden 5.

\section{Filtro pasabajas de orden 5}

\subsection{Frecuencia de corte escogida}

La senoidal debe conservar la fundamental de la onda triangular y atenuar sus
armonicos. La fundamental esta en:

\[
f_{tri}=\frac{f_0}{2}.
\]

El primer armonico no deseado de la triangular es el tercero:

\[
f_{3}=3f_{tri}.
\]

Se escoge:

\[
f_c=1.25f_{tri}.
\]

Para el ejemplo:

\[
f_c=1.25(615)=768.75~\text{Hz}.
\]

La relacion del tercer armonico respecto a \(f_c\) es:

\[
\frac{f_3}{f_c}
=
\frac{3f_{tri}}{1.25f_{tri}}
=
\frac{3}{1.25}
=2.4.
\]

Entonces el tercer armonico queda a \(2.4f_c\), donde un filtro de quinto orden
ya proporciona atenuacion apreciable.

\subsection{Tipo de respuesta segun Z}

El enunciado asigna el tipo de respuesta con el digito \(Z\):

\[
\begin{array}{c|c}
\text{Valor de } Z & \text{Respuesta} \\
\hline
1,2,3 & \text{Bessel} \\
4,5,6 & \text{Butterworth} \\
7,8,9 & \text{Chebyshev}
\end{array}
\]

Para el ejemplo:

\[
Z=3.
\]

Entonces corresponde:

\[
\text{Respuesta Bessel}.
\]

\subsection{Descomposicion en etapas}

Un filtro de orden 5 se implementa como:

\[
5 = 1 + 2 + 2.
\]

Por tanto se usan:

\begin{itemize}
  \item una etapa de primer orden,
  \item una etapa Sallen-Key de segundo orden,
  \item otra etapa Sallen-Key de segundo orden.
\end{itemize}

La etapa de primer orden tiene:

\[
f_p=\frac{1}{2\pi RC}.
\]

Se despeja:

\[
R=\frac{1}{2\pi f_p C}.
\]

Cada etapa Sallen-Key se disena con:

\[
R_1=R_2=R,
\qquad
C_1=C_2=C.
\]

Para esa seleccion:

\[
f_n=\frac{1}{2\pi RC}.
\]

Se despeja:

\[
R=\frac{1}{2\pi f_n C}.
\]

La ganancia no inversora de cada Sallen-Key es:

\[
K=1+\frac{R_f}{R_g}.
\]

Con componentes iguales, la relacion de calidad usada es:

\[
Q=\frac{1}{3-K}.
\]

Se despeja \(K\):

\[
Q(3-K)=1.
\]

\[
3Q-KQ=1.
\]

\[
-KQ=1-3Q.
\]

\[
K=\frac{3Q-1}{Q}.
\]

\[
K=3-\frac{1}{Q}.
\]

\subsection{Coeficientes usados}

\[
\begin{array}{c|c|c|c|c|c|c|c}
\text{Respuesta} & f_{p1}/f_c & f_{n1}/f_c & Q_1 & K_1 &
f_{n2}/f_c & Q_2 & K_2 \\
\hline
\text{Bessel} & 1.5023 & 1.5563 & 0.5635 & 1.2254
& 1.7554 & 0.9165 & 1.9089 \\
\text{Butterworth} & 1.0000 & 1.0000 & 0.6180 & 1.3820
& 1.0000 & 1.6180 & 2.3820 \\
\text{Chebyshev 0.5 dB} & 0.3623 & 0.6905 & 1.1778 & 2.1511
& 1.0177 & 4.5450 & 2.7800
\end{array}
\]

Para \(Z=3\), se toma la fila Bessel:

\[
f_{p1}=1.5023f_c,
\qquad
f_{n1}=1.5563f_c,
\qquad
f_{n2}=1.7554f_c.
\]

\subsection{Calculo de la etapa de primer orden}

Se escogio:

\[
C_f=10~\text{nF}=10\times 10^{-9}~\text{F}.
\]

Primero:

\[
f_{p1}=1.5023(768.75)=1155.5~\text{Hz}.
\]

La resistencia es:

\[
R_{p1}=
\frac{1}{2\pi f_{p1}C_f}.
\]

Sustituyendo:

\[
R_{p1}
=
\frac{1}{2\pi(1155.5)(10\times 10^{-9})}.
\]

\[
R_{p1}=13780.8~\Omega.
\]

Por tanto:

\[
R_{p1}\approx 13.8~\text{k}\Omega.
\]

\subsection{Calculo de la primera etapa Sallen-Key}

Para la primera etapa Bessel:

\[
Q_1=0.5635.
\]

La ganancia requerida es:

\[
K_1=3-\frac{1}{Q_1}.
\]

Sustituyendo:

\[
K_1=3-\frac{1}{0.5635}.
\]

\[
K_1=3-1.7745=1.2255.
\]

La frecuencia natural es:

\[
f_{n1}=1.5563f_c.
\]

\[
f_{n1}=1.5563(768.75)=1196.0~\text{Hz}.
\]

La resistencia de la etapa es:

\[
R_{SK1}
=
\frac{1}{2\pi f_{n1}C_f}.
\]

Sustituyendo:

\[
R_{SK1}
=
\frac{1}{2\pi(1196.0)(10\times 10^{-9})}
=13302.4~\Omega.
\]

Por tanto:

\[
R_{SK1}\approx 13.3~\text{k}\Omega.
\]

Para implementar la ganancia se usa:

\[
K_1=1+\frac{R_{f1}}{R_{g1}}.
\]

Si se escoge:

\[
R_{g1}=10~\text{k}\Omega,
\]

entonces:

\[
R_{f1}=(K_1-1)R_{g1}.
\]

\[
R_{f1}=(1.2255-1)(10~\text{k}\Omega)
=2.255~\text{k}\Omega.
\]

\subsection{Calculo de la segunda etapa Sallen-Key}

Para la segunda etapa Bessel:

\[
Q_2=0.9165.
\]

La ganancia requerida es:

\[
K_2=3-\frac{1}{Q_2}.
\]

Sustituyendo:

\[
K_2=3-\frac{1}{0.9165}.
\]

\[
K_2=3-1.0911=1.9089.
\]

La frecuencia natural es:

\[
f_{n2}=1.7554f_c.
\]

\[
f_{n2}=1.7554(768.75)=1349.7~\text{Hz}.
\]

La resistencia de la etapa es:

\[
R_{SK2}
=
\frac{1}{2\pi f_{n2}C_f}.
\]

Sustituyendo:

\[
R_{SK2}
=
\frac{1}{2\pi(1349.7)(10\times 10^{-9})}
=11794.1~\Omega.
\]

Por tanto:

\[
R_{SK2}\approx 11.8~\text{k}\Omega.
\]

Para implementar la ganancia:

\[
K_2=1+\frac{R_{f2}}{R_{g2}}.
\]

Si:

\[
R_{g2}=10~\text{k}\Omega,
\]

entonces:

\[
R_{f2}=(K_2-1)R_{g2}.
\]

\[
R_{f2}=(1.9089-1)(10~\text{k}\Omega)
=9.089~\text{k}\Omega.
\]

\subsection{Atenuador de entrada del filtro}

Las dos etapas Sallen-Key tienen ganancia en banda pasante. La ganancia total de
esas dos etapas es:

\[
K_T=K_1K_2.
\]

Sustituyendo:

\[
K_T=(1.2255)(1.9089)=2.3393.
\]

Para que el filtro completo no amplifique demasiado la triangular antes de
filtrar, se coloca un atenuador resistivo al inicio con ganancia:

\[
A=\frac{1}{K_T}.
\]

El divisor resistivo cumple:

\[
A=\frac{R_{bot}}{R_{top}+R_{bot}}.
\]

Se despeja \(R_{top}\):

\[
A(R_{top}+R_{bot})=R_{bot}.
\]

\[
AR_{top}+AR_{bot}=R_{bot}.
\]

\[
AR_{top}=R_{bot}-AR_{bot}.
\]

\[
R_{top}=\frac{R_{bot}(1-A)}{A}.
\]

Como \(A=1/K_T\):

\[
R_{top}=R_{bot}(K_T-1).
\]

Se escogio:

\[
R_{bot}=100~\text{k}\Omega.
\]

Entonces:

\[
R_{top}
=
(100~\text{k}\Omega)(2.3393-1)
=133.93~\text{k}\Omega.
\]

Por tanto:

\[
R_{top}\approx 134~\text{k}\Omega.
\]

\section{Resumen numerico para el ejemplo \(X=1,Y=2,Z=3\)}

\[
\begin{array}{c|c}
\text{Parametro} & \text{Valor} \\
\hline
f_0 & 1230~\text{Hz} \\
f_{saw} & 2460~\text{Hz} \\
f_{tri} & 615~\text{Hz} \\
f_c & 768.75~\text{Hz} \\
R_{osc,LTspice} & 38.7~\text{k}\Omega \\
I_{saw} & 123~\mu\text{A} \\
R_{tri} & 20.3~\text{k}\Omega \\
R_{p1} & 13.8~\text{k}\Omega \\
R_{SK1} & 13.3~\text{k}\Omega \\
R_{SK2} & 11.8~\text{k}\Omega \\
R_{top} & 134~\text{k}\Omega \\
R_{bot} & 100~\text{k}\Omega
\end{array}
\]

\end{document}
